\documentclass[10pt,a4paper,fleqn]{article}
\usepackage[margin=0.55in]{geometry}
\usepackage{graphicx,float,amsmath,amssymb,booktabs,enumitem,multicol,titlesec,fancyhdr,tikz}
\setlength{\parskip}{16pt}
\setlength{\parindent}{0pt}
\linespread{1.7}
\titleformat{\section}{\Huge\bfseries}{\thesection}{0.2em}{}
\titleformat{\subsection}{\Large\bfseries}{\thesubsection}{0.1em}{}
\pagestyle{fancy}\fancyhf{}
\lhead{\tiny A Very Long Running Header That Carries No Information Whatsoever And Collides With The Rule}
\rhead{\tiny \today}
\title{\Huge A Title So Large That It Eats The Entire First Page Of The Manuscript}
\author{A. Student \and B. Student \and C. Student \and D. Student}
\date{}
\begin{document}
\maketitle
This opening paragraph is long enough to create an awkward block right under the
title. It repeats itself, spaces badly and mixes several poor decisions. The
point is not to communicate a result but to stress every mechanism at once:
headings, equations, figures, tables, lists, margins and page breaks.
\section{Introduction}
Layout optimization for academic manuscripts is a constrained problem: the visual result of a document is a global property of the whole file, while most editing actions are local. A change that removes one bad line can push a float to the next page and create a much larger gap, so an optimizer must always recompile and re-evaluate the entire document before accepting an edit.

The scoring model used here separates logic from aesthetics. Logic asks whether the file compiles, whether the rendered PDF keeps the body text byte-identical, and whether the requested hard constraints such as font size, margins and page limits are satisfied. Aesthetics collects measurable layout defects such as overfull lines, underfull lines, badly placed floats and inconsistent vertical spacing.

A long unbreakable token follows: supercalifragilisticexpialidociousmicroscopicsilicovolcanoconiosispneumonoultramicroscopicsilicovolcanoconiosis.

\noindent Manually de-indented paragraph one.
\noindent Manually de-indented paragraph two.
\noindent Manually de-indented paragraph three.

\vspace{3cm}
\vspace{-1.5cm}

Emphasis is all over the place: \underline{one}, \underline{two},
\underline{three}, \underline{four} and {\Large a size switch} inside the
running text, plus {\tiny a tiny fragment}.

Bare dollar math appears four times:
$$a^2 + b^2 = c^2$$
$$e^{i\pi} + 1 = 0$$
$$\frac{\partial u}{\partial t} = \alpha \nabla^2 u$$
$$\int_0^1 x^2\,dx = \frac{1}{3}$$

\newpage


\section{Mathematical Model}
The pipeline is deterministic. Given the same source file, the same requirement specification and the same toolchain, it produces the same trace of candidate actions, the same accept and rollback decisions, and the same final document. Determinism matters for reproducibility, for regression testing and for the audit trail that is written next to every run.

Perception is layered. The source layer parses the LaTeX file and looks for patterns that are known to be unstable, such as manual page breaks, manual vertical spacing, bare dollar math and float specifications that pin a float to its declaration point. The log layer reads the compiler log and counts overfull boxes, underfull boxes and badness warnings.

The rendered layer measures the PDF itself. Page images give coverage statistics, the vertical position of ink, and the size of the largest empty band on each page. Those numbers are coarse on purpose: they are used as evidence that something looks wrong, not as a model of human taste.

Float placement is the hardest part of the problem. A figure that is referenced in one paragraph may be typeset several pages later, and the quality of the page depends on whether the top, bottom and text fraction constraints of the float mechanism are balanced. Sanitising the float specification to the top, bottom and page forms is a cheap first step.

\begin{center}
This whole paragraph of running text is centred, which destroys the left edge
that the reader uses to navigate the page.
\end{center}
\begin{flushright}
And this one is flushed right for no reason at all.
\end{flushright}

\begin{align}
y &= a_0 + a_1 x + a_2 x^2 + a_3 x^3 + a_4 x^4 + a_5 x^5 + a_6 x^6 + a_7 x^7 \\
z &= \frac{\alpha + \beta + \gamma + \delta + \epsilon}{1 + \alpha^2 + \beta^2 + \gamma^2}
\end{align}

\begin{figure}[H]\centering
\includegraphics[width=1.25\linewidth]{example-image}
\caption{A figure that is much wider than the text block.}
\end{figure}
\begin{figure}[h]\centering
\includegraphics[width=0.9\linewidth,height=0.55\textheight,keepaspectratio]{example-image}
\caption{A figure that eats more than half of the page height.}
\end{figure}

\begin{table}[!ht]\centering
\caption{A dense table with mixed rules.}
\begin{tabular}{lrrrrrrrr}\hline
Method & Accuracy & Precision & Recall & F1 & Time & Memory & Score & Rank\\\hline
Method A & 0.912345 & 0.934567 & 0.901234 & 0.917321 & 123.45 & 1024 & 98.123 & 1 \\\toprule
Method B & 0.901234 & 0.923456 & 0.889012 & 0.905432 & 234.56 & 2048 & 91.456 & 2 \\\bottomrule
\end{tabular}
\end{table}
\begin{table}[!ht]\centering
\caption{Another table with narrow fixed columns.}
\begin{tabular}{p{2.0cm}p{2.0cm}p{2.0cm}}
\hline Very long heading that wraps & Another long heading & Third heading \\\hline
A large amount of text in a cell & More verbose text & Additional information \\\hline
\end{tabular}
\end{table}

\begin{itemize}[leftmargin=1pt,itemsep=18pt,topsep=16pt]
\item A list with far too much vertical spacing between the items.
\item A second item, followed by a nested list:
\begin{enumerate}[itemsep=12pt]\item nested one\item nested two\end{enumerate}
\end{itemize}

\begin{multicols}{2}
This mid-document two-column region has no reason to exist. The columns are
narrow, hyphenation explodes and the page rhythm is broken.
\subsection{A Small Subsection} Short text.
\end{multicols}

An extreme closing remark: \hfill \vfill

A footnote is abused here\footnote{This footnote is just as long as the
paragraph it annotates, which is not how footnotes work.}\footnote{A second
footnote on the very same line.} and the sentence keeps going for a while so
that the page has to stretch.

A rotate box is used for no reason: \rotatebox{90}{rotated text}.

\clearpage


\section{Experiments}
Intervention cost makes the optimizer conservative. Every edit that is applied to the source increases the intervention score, so a repair is only kept when the total objective improves after a full recompilation. This is what prevents the system from grooming its own score with edits that no human would ever make.

Page limits deserve special care. Compressing a document by shrinking its margins changes the relationship between floats and the text that refers to them, and the resulting layout often looks wrong even though every number in the report is legal. For that reason the default policy refuses to touch the global text block and reports an honest failure instead.

The pipeline is deterministic. Given the same source file, the same requirement specification and the same toolchain, it produces the same trace of candidate actions, the same accept and rollback decisions, and the same final document. Determinism matters for reproducibility, for regression testing and for the audit trail that is written next to every run.

Perception is layered. The source layer parses the LaTeX file and looks for patterns that are known to be unstable, such as manual page breaks, manual vertical spacing, bare dollar math and float specifications that pin a float to its declaration point. The log layer reads the compiler log and counts overfull boxes, underfull boxes and badness warnings.

The rendered layer measures the PDF itself. Page images give coverage statistics, the vertical position of ink, and the size of the largest empty band on each page. Those numbers are coarse on purpose: they are used as evidence that something looks wrong, not as a model of human taste.

A duplicate label lives here \label{sec:dup} and the same label was
already used earlier in the document.

\pagebreak


\section{Conclusion}
The rendered layer measures the PDF itself. Page images give coverage statistics, the vertical position of ink, and the size of the largest empty band on each page. Those numbers are coarse on purpose: they are used as evidence that something looks wrong, not as a model of human taste.

Float placement is the hardest part of the problem. A figure that is referenced in one paragraph may be typeset several pages later, and the quality of the page depends on whether the top, bottom and text fraction constraints of the float mechanism are balanced. Sanitising the float specification to the top, bottom and page forms is a cheap first step.


\section*{Appendix}
Text in the appendix.
\end{document}
